% GENERATED FILE. Edit canonical lessons, then run npm run generate:books.
% canonical-source-hash: fnv1a64:0cb198b687947e9d
% canonical-lessons: ZH-C122-yingye, ZH-C122-kaishi, ZH-C122-jieshu, ZH-C122-shikebiao, ZH-C122-dazhe, ZH-R122-review-117-122, ZH-R122-review-final-122

\chapter{Opening Hours}
\label{ch:zh-opening-hours}
\hlchaptermodality{122}

\begin{chapteropening}
\textbf{By the end of this chapter:} \emph{I can say open for business, start, finish, discount and timetable.}

Complete the last lesson of chapter 122: I can say open for business, start, finish, discount and timetable.
\end{chapteropening}

\section[yíngyè]{yíngyè — to be open for business}

\label{lesson:ZH-C122-yingye}



\begin{quote}
Before the new one: say the Mandarin for to revise, then the Mandarin for wrong.
\end{quote}

\subsection*{You'll want to know: yíngyè}
\emph{yíngyè} — \textquotedblleft{}to be open for business\textquotedblright{}.

\begin{quote}
\emph{shāngdiàn jǐ diǎn yíngyè} — What time does the shop open?
\end{quote}

Its characters come later; for now it is a word for the ear.

\begin{cousinweb}
\emph{yíngyè} is built from yíng and yè, both inherited from Middle Chinese, \emph{yweng} and \emph{ngjaep} (Guangyun, in Baxter's transcription), glossed 'to build, to plan' and 'affairs'.
\end{cousinweb}

\begin{grammarlens}[title={What you've built}]
One more word for signs and notices.
\end{grammarlens}

\subsection*{Your turn}
\begin{itemize}
\raggedright
  \item \emph{Say it:} \emph{yíngyè}
  \item \emph{Say it:} \emph{yíngyè}, once more
  \item \emph{Say it:} \emph{shāngdiàn jǐ diǎn yíngyè} — What time does the shop open?
  \item \emph{Recall:} say the Mandarin for to revise, then the Mandarin for wrong, then say \emph{yíngyè} again
\end{itemize}

\begin{tcolorbox}[breakable,colback=teal!4,colframe=teal!35!black,title={Before you move on}]
What does yíngyè mean? (\textquotedblleft{}To be open for business\textquotedblright{}.) Say it once more.
\end{tcolorbox}
\section[kāishǐ]{kāishǐ — to begin; to start}

\label{lesson:ZH-C122-kaishi}



\begin{quote}
Before the new one: say the Mandarin for wrong, then the Mandarin for to be open for business.
\end{quote}

\subsection*{You'll want to know: kāishǐ}
\emph{kāishǐ} — \textquotedblleft{}to begin; to start\textquotedblright{}.

\begin{quote}
\emph{wǒmen jǐ diǎn kāishǐ} — What time do we start?
\end{quote}

Its characters come later; for now it is a word for the ear.

\begin{cousinweb}
\emph{kāishǐ} is built from \zh{开} 'to open' and shǐ 'beginning', both inherited from Middle Chinese, \emph{khoj} and \emph{syiX} (Guangyun, in Baxter's transcription).
\end{cousinweb}

\begin{grammarlens}[title={What you've built}]
One more word for signs and notices.
\end{grammarlens}

\subsection*{Your turn}
\begin{itemize}
\raggedright
  \item \emph{Say it:} \emph{kāishǐ}
  \item \emph{Say it:} \emph{kāishǐ}, once more
  \item \emph{Say it:} \emph{wǒmen jǐ diǎn kāishǐ} — What time do we start?
  \item \emph{Recall:} say the Mandarin for wrong, then the Mandarin for to be open for business, then say \emph{kāishǐ} again
\end{itemize}

\begin{tcolorbox}[breakable,colback=teal!4,colframe=teal!35!black,title={Before you move on}]
What does kāishǐ mean? (\textquotedblleft{}To begin; to start\textquotedblright{}.) Say it once more.
\end{tcolorbox}
\section[jiéshù]{jiéshù — to come to an end; to finish}

\label{lesson:ZH-C122-jieshu}



\begin{quote}
Before the new one: say the Mandarin for to be open for business, then the Mandarin for to begin.
\end{quote}

\subsection*{You'll want to know: jiéshù}
\emph{jiéshù} — \textquotedblleft{}to come to an end; to finish\textquotedblright{}.

\begin{quote}
\emph{kè jǐ diǎn jiéshù} — What time does the lesson finish?
\end{quote}

Its characters come later; for now it is a word for the ear.

\begin{cousinweb}
\emph{jiéshù} is built from jié 'to tie' and shù 'to bind', both inherited from Middle Chinese, \emph{ket} and \emph{syowk} (Guangyun, in Baxter's transcription).
\end{cousinweb}

\begin{grammarlens}[title={What you've built}]
One more word for signs and notices.
\end{grammarlens}

\subsection*{Your turn}
\begin{itemize}
\raggedright
  \item \emph{Say it:} \emph{jiéshù}
  \item \emph{Say it:} \emph{jiéshù}, once more
  \item \emph{Say it:} \emph{kè jǐ diǎn jiéshù} — What time does the lesson finish?
  \item \emph{Recall:} say the Mandarin for to be open for business, then the Mandarin for to begin, then say \emph{jiéshù} again
\end{itemize}

\begin{tcolorbox}[breakable,colback=teal!4,colframe=teal!35!black,title={Before you move on}]
What does jiéshù mean? (\textquotedblleft{}To come to an end; to finish\textquotedblright{}.) Say it once more.
\end{tcolorbox}
\section[shíkèbiǎo]{shíkèbiǎo — timetable}

\label{lesson:ZH-C122-shikebiao}



\begin{quote}
Before the new one: say the Mandarin for to begin, then the Mandarin for to come to an end.
\end{quote}

\subsection*{You'll want to know: shíkèbiǎo}
\emph{shíkèbiǎo} — \textquotedblleft{}timetable\textquotedblright{}.

\begin{quote}
\emph{wǒ yǒu yì zhāng huǒchē shíkèbiǎo} — I've got a train timetable.
\end{quote}

Its characters come later; for now it is a word for the ear.

\begin{cousinweb}
\emph{shíkèbiǎo} is built from \zh{时} 'time', kè 'a notch, a quarter of an hour' and biǎo 'table', all inherited from Middle Chinese, \emph{dzyi}, \emph{khok} and \emph{pjewX} (Guangyun, in Baxter's transcription); the dictionary glosses the second 'to carve' and the third 'to make clear'.
\end{cousinweb}

\begin{grammarlens}[title={What you've built}]
One more word for signs and notices.
\end{grammarlens}

\subsection*{Your turn}
\begin{itemize}
\raggedright
  \item \emph{Say it:} \emph{shíkèbiǎo}
  \item \emph{Say it:} \emph{shíkèbiǎo}, once more
  \item \emph{Say it:} \emph{wǒ yǒu yì zhāng huǒchē shíkèbiǎo} — I've got a train timetable.
  \item \emph{Recall:} say the Mandarin for to begin, then the Mandarin for to come to an end, then say \emph{shíkèbiǎo} again
\end{itemize}

\begin{tcolorbox}[breakable,colback=teal!4,colframe=teal!35!black,title={Before you move on}]
What does shíkèbiǎo mean? (\textquotedblleft{}Timetable\textquotedblright{}.) Say it once more.
\end{tcolorbox}
\section[dǎzhé]{dǎzhé — to give a discount; to be on sale}

\label{lesson:ZH-C122-dazhe}



\begin{quote}
Before the new one: say the Mandarin for to come to an end, then the Mandarin for timetable.
\end{quote}

\subsection*{You'll want to know: dǎzhé}
\emph{dǎzhé} — \textquotedblleft{}to give a discount; to be on sale\textquotedblright{}.

\begin{quote}
\emph{zhège xīngqī dǎzhé} — Things are on sale this week.
\end{quote}

Its characters come later; for now it is a word for the ear.

\begin{cousinweb}
\emph{dǎzhé} is built from \zh{打} 'to strike, to do' and zhé, both inherited from Middle Chinese, \emph{taengX} and \emph{tsyet} (Guangyun, in Baxter's transcription); the dictionary glosses the second 'to bend and break'. How the pair came to mean a discount is not settled.
\end{cousinweb}

\begin{grammarlens}[title={What you've built}]
One more word for signs and notices.
\end{grammarlens}

\subsection*{Your turn}
\begin{itemize}
\raggedright
  \item \emph{Say it:} \emph{dǎzhé}
  \item \emph{Say it:} \emph{dǎzhé}, once more
  \item \emph{Say it:} \emph{zhège xīngqī dǎzhé} — Things are on sale this week.
  \item \emph{Recall:} say the Mandarin for to come to an end, then the Mandarin for timetable, then say \emph{dǎzhé} again
\end{itemize}

\begin{tcolorbox}[breakable,colback=teal!4,colframe=teal!35!black,title={Before you move on}]
What does dǎzhé mean? (\textquotedblleft{}To give a discount; to be on sale\textquotedblright{}.) Say it once more.
\end{tcolorbox}
\section[(dialogue)]{Looking back: from Which Way? to Opening Hours}

\label{lesson:ZH-R122-review-117-122}



\begin{quote}
Say the Mandarin for the east and the south.
\end{quote}

\subsection*{Your turn}
\begin{itemize}
\raggedright
  \item \emph{Say it:} the east, the south, the west, the north, direction
  \item \emph{Say it:} parcel, to post, letter, text message, to receive
  \item \emph{Say it:} schoolbag, suitcase, small box, carrier bag, luggage
  \item \emph{Say it:} takeaway, to make a booking, address, to pack food to take away, dim sum
  \item \emph{Say it:} to teach, to practise, to answer, to revise, wrong
  \item \emph{Say it:} to be open for business, to begin, to come to an end, timetable, to give a discount
  \item \emph{Read it:} \zh{包} — each one aloud, then a word that holds it
\end{itemize}

\begin{tcolorbox}[breakable,colback=teal!4,colframe=teal!35!black,title={Before you move on}]
The east? (\emph{dōngbian}.) The south? (\emph{nánbian}.) The west? (\emph{xībian}.) The north? (\emph{běibian}.) Direction? (\emph{fāngxiàng}.) Parcel? (\emph{bāoguǒ}.) To post? (\emph{jì}.) Letter? (\emph{xìn}.) Text message? (\emph{duǎnxìn}.) To receive? (\emph{shōu}.) Schoolbag? (\textbf{\zh{书包}}, \emph{shūbāo}.) Suitcase? (\emph{xiāngzi}.) Small box? (\emph{hézi}.) Carrier bag? (\emph{dàizi}.) Luggage? (\emph{xíngli}.) Takeaway? (\emph{wàimài}.) To make a booking? (\emph{dìng}.) Address? (\emph{dìzhǐ}.) To pack food to take away? (\textbf{\zh{打包}}, \emph{dǎbāo}.) Dim sum? (\textbf{\zh{点心}}, \emph{diǎnxin}.) To teach? (\emph{jiāo}.) To practise? (\emph{liànxí}.) To answer? (\emph{huídá}.) To revise? (\emph{fùxí}.) Wrong? (\emph{cuò}.) To be open for business? (\emph{yíngyè}.) To begin? (\emph{kāishǐ}.) To come to an end? (\emph{jiéshù}.) Timetable? (\emph{shíkèbiǎo}.) To give a discount? (\emph{dǎzhé}.)
\end{tcolorbox}
\section[(dialogue)]{Opening Hours — once more}

\label{lesson:ZH-R122-review-final-122}



\begin{quote}
Say the Mandarin for to be open for business and to begin.
\end{quote}

\subsection*{Your turn}
\begin{itemize}
\raggedright
  \item \emph{Say it:} to be open for business, to begin, to come to an end, timetable, to give a discount
\end{itemize}

\begin{tcolorbox}[breakable,colback=teal!4,colframe=teal!35!black,title={Before you move on}]
To be open for business? (\emph{yíngyè}.) To begin? (\emph{kāishǐ}.) To come to an end? (\emph{jiéshù}.) Timetable? (\emph{shíkèbiǎo}.) To give a discount? (\emph{dǎzhé}.)
\end{tcolorbox}
